\documentclass[10pt,a4paper]{article}
\usepackage[utf8]{inputenc}
\usepackage{booktabs}
\usepackage[a4paper,left=3cm,right=2cm]{geometry}
\usepackage{amsmath}
\usepackage[thmmarks, amsmath, thref]{ntheorem}
\usepackage{amsfonts}
\usepackage{amssymb}
\usepackage{fancyhdr}
\pagestyle{fancy}
\usepackage{anysize}
\usepackage{graphicx}
\usepackage{float}
\usepackage{hyperref}
\usepackage{enumitem}
\usepackage{verbatim}
\newtheorem{theorem}{Theorem}
\theoremsymbol{\ensuremath{\blacksquare}}
\newtheorem*{solution}{Solución:}

\usepackage{xcolor}
\usepackage{listings}
% Margenes y formalidades

\textwidth 6.25in % Margen Derecho y ancho del texto
\textheight 8.5in
\oddsidemargin 0in % Margen Izquierdo
\headheight 0in % Largo del texto
\leftmargin 1in
\parindent 0em % Salto entre la indentación (Sangría)https://www.overleaf.com/4519127669sdpptkhvsbfn
\topmargin -0.5in % Margen Superior
\parskip 0.5ex % Salto entre parrafos
\baselineskip 1pt



\lhead{
\includegraphics[scale=0.08]{Figuras/logoUSM.png}} % texto izquierda de la cabecera
%\chead{}                                      % texto centro de la cabecera
\rhead{Universidad Técnica Federico Santa María\\Ingeniería Comercial } % número de página a la derecha
\renewcommand{\headrulewidth}{0.4pt} % grosor de la línea de la cabecera
\renewcommand{\footrulewidth}{0.4pt} % grosor de la línea del pie

\begin{document}
\vspace{5cm}



\begin{center}
\begin{large}
\textbf{\\CERTAMEN 3:}\textbf{ ECONOMETRÍA ICS-013}
\end{large}\\
\vspace{0.2cm}
II Semestre 2025\\
Profesores: Franco Antonucci Fernández\\
%Ayudante: 
\end{center}

\textbf{Instrucciones:} 

\begin{itemize}

    \item Usted dispone de 90 minutos  para intentar los 100 puntos del certamen.
    \item Cuide la presentación y redacción de sus preguntas.
    \item No puede utilizar celular, computador o tablet, pero sí una calculadora TI.
     
\end{itemize}

\section{RdeA9.1 Comentes y preguntas cortas (40 puntos)}
{\bf Elija y responda solo 4 de las siguientes preguntas.}
\begin{enumerate}
    \item Considere la base de datos \texttt{bwght} de Wooldridge, que contiene información sobre recién nacidos
y características de sus madres durante el embarazo. Se define la variable dependiente
$lowbirth_i$ como una variable dicotómica que toma valor 1 si el bebé $i$ nace con bajo peso
(por ejemplo, menos de 2500 gramos), y 0 en caso contrario.

A partir de la base, se ha estimado el siguiente modelo de probabilidad lineal (MPL):

\[
\widehat{lowbirth}_i
= 0.20 + 0.02\, cigs_i - 0.004\, faminc_i,
\]

donde:
\begin{itemize}
    \item $cigs_i$: número de cigarrillos que la madre fuma por día durante el embarazo,
    \item $faminc_i$: ingreso familiar medido en miles de dólares anuales.
\end{itemize}

Responda:

\begin{enumerate}
    \item (2 puntos) Interprete el coeficiente $\hat\beta_0 = 0.20$.  
    \item (2 puntos) Interprete el coeficiente asociado a la variable $cigs_i$.
    \item (2 puntos) Interprete el coeficiente asociado a la variable $faminc_i$.
    \item (2 puntos)  Calcule la probabilidad de bajo peso para un bebé cuya madre fuma 10 cigarrillos al día y tiene ingreso familiar de 20 (miles de dólares).
    \item (2 puntos)  ¿Qué problema del modelo de probabilidad lineal se evidencia al considerar valores extremos de las variables?


\end{enumerate}
    \item Luego de estimar un modelo de regresión, se aplica el test de Breusch--Pagan para evaluar la presencia de heterocedasticidad en los residuos. El test arroja un \textit{p-valor} de 0.004.

\begin{enumerate}
  \item[a)] ¿Qué implica este resultado respecto al cumplimiento del supuesto de homocedasticidad en el modelo estimado?
  \item[b)] ¿Qué efectos tiene la heterocedasticidad sobre los errores estándar y sobre la validez de la inferencia estadística basada en el modelo?
  \item[c)] ¿Qué métodos pueden utilizarse para corregir o mitigar este problema?
\end{enumerate}

    \item En un estudio que relaciona el promedio de calificaciones en la universidad con el tiempo dedicado a diversas actividades, usted reparte un cuestionario a varios estudiantes en el que les pregunta cuántas horas destinan semanalmente a estas cuatro categorías: estudiar, dormir, trabajar y tiempo libre. Cualquier actividad se sitúa en una de las categorías, de modo que la suma de horas de cada estudiante en las cuatro debe ser 168.

\begin{enumerate}
    \item[(i)] (5 puntos) En el modelo, explique por qué el modelo infringe la suposición de no Colinealidad Perfecta.

    \[
        GPA = \beta_0 + \beta_1 \, estud + \beta_2 \, dorm + \beta_3 \, trab + \beta_4 \, tlib + u,
    \]

    \item[(ii)] (5 puntos) ¿Cómo reformularía el modelo para que sus parámetros tuvieran una interpretación  útil y satisficieran la suposición de Colinealidad Imperfecta?

\end{enumerate}

    
    \item Si en un modelo existen problemas por variable omitida y este se puede solucionar con una variable instrumental, no necesito saber cuál es la variable omitida, sólo que existe una. Comente.\\

    
    \item Un investigador ha simulado tres series temporales de 200 observaciones. 
Cada una corresponde a uno de los siguientes procesos:

\begin{itemize}
    \item Un proceso de \textbf{ruido blanco} (white noise).
    \item Un proceso \textbf{MA(1)} con parámetro $\theta = 0.8$:
    \[
        X_t = \varepsilon_t + 0.8\,\varepsilon_{t-1},
    \]
    \item Un proceso \textbf{AR(1)} con parámetro $\phi = 0.8$:
    \[
        X_t = 0.8\,X_{t-1} + \varepsilon_t,
    \]
\end{itemize}

donde en todos los casos $\varepsilon_t \sim \text{i.i.d. } N(0,1)$.

Para cada caso se presentan dos gráficos:
\begin{itemize}
    \item A la izquierda: la serie temporal simulada $\{X_t\}$.
    \item A la derecha: su función de autocorrelación (ACF) hasta el rezago 30, con bandas de confianza aproximadas.
\end{itemize}

Los gráficos han sido rotulados como \textbf{Caso A}, \textbf{Caso B} y \textbf{Caso C} (sin indicar el tipo de proceso que representan).

\begin{figure}[H]
    \centering
    \includegraphics[width=0.65\linewidth]{certamenes//certamenesDiurno25-2/seriesACF.png}
\end{figure}



\begin{enumerate}
    \item[(a)] (4 puntos) Identifique cuál de los tres casos corresponde al proceso de \textbf{ruido blanco}. Justifique su respuesta con base en el comportamiento de la serie y en la forma de su ACF.

    \item[(b)] (6 puntos) Para los otros dos casos, indique cuál corresponde al proceso \textbf{MA(1)} con $\theta = 0.8$ y cuál corresponde al proceso \textbf{AR(1)} con $\phi = 0.8$. Justifique su elección explicando cómo se refleja en la ACF el hecho de que:
    \begin{itemize}
        \item En el MA(1) la autocorrelación se “corta” después de un número finito de rezagos.
        \item En el AR(1) la autocorrelación decae gradualmente de forma geométrica.
    \end{itemize}
\end{enumerate}

    \item En la presentación de Ph.D. Cristóbal Otero de su investigación \textit{Physicians’ Occupational Licensing and the Quantity-Quality trade-off} se distinguen dos efectos causales, en acceso a la salud y en mortalidad. ¿Qué estrategia empírica utilizan los autores para determinar la causalidad? ¿Cómo son esos efectos en términos de significancia y dirección?\\

\end{enumerate}

\section{RdeA9.1 Ejercicios (60 puntos)}

\begin{enumerate}
    \item (20 puntos) Considere el siguiente modelo que busca explicar el impacto innovador de un proyecto tecnológico,
medido mediante el índice \textit{ImpactScore}:

\[
ImpactScore_i
= \beta_0 
+ \beta_1\, Invest_i
+ \beta_2\, TeamExp_i
+ \beta_3\, Collab_i
+ u_i,
\]

donde:
\begin{itemize}
    \item $Invest_i$: inversión inicial destinada al proyecto (en miles de dólares),
    \item $TeamExp_i$: años totales de experiencia acumulada por el equipo fundador,
    \item $Collab_i$: número de alianzas con universidades, laboratorios u otras empresas.
\end{itemize}

La organización desea investigar si la relación entre estas variables es igual para dos tipos de startups:
\textit{startups verdes} (grupo 1) y \textit{startups de inteligencia artificial} (grupo 2). Para ello realiza un Test de Chow.

Los resultados obtenidos son:
\[
SSR_1 = 210, \qquad SSR_2 = 250, \qquad SSR_r = 500,
\]
con $n = 120$ observaciones totales.

Responda:

\begin{enumerate}
    \item[(a)] (5 puntos) Formule las hipótesis nula y alternativa del Test de Chow.

    \item[(b)] (10 puntos) Indique si debe rechazarse la hipótesis nula. (considere el valor crítico al 5\%,
    $F_{0.05}(4,112)=2.45$.
    
    \item[(c)] (5 puntos) Interprete el resultado en el contexto de la innovación tecnológica.
\end{enumerate}

    \item (20 puntos) Se dispone de la siguiente serie del PIB real de Chile, medida en millones de pesos a precios constantes, para el período 1999--2024:

\begin{center}
    \begin{figure}[H]
        \centering
        \includegraphics[width=0.7\linewidth]{certamenes/Certamenes T2/imagen pib.png}
        \caption{PIB real de Chile (1999--2024)}
        \label{fig:pib_chile}
    \end{figure}
\end{center}

Un investigador está interesado en analizar las propiedades dinámicas del PIB real. 
\begin{enumerate}
    \item[(a)]  (8 puntos) Suponiendo que el investigador plantea que el PIB real podría ser una serie integrada de orden 1, escriba el modelo de regresión auxiliar del test Dickey--Fuller  correspondiente. Indique claramente la forma funcional, la variable dependiente y el término cuyo coeficiente es relevante para el contraste de raíz unitaria.

    \item[(c)] (4 puntos) Formule la hipótesis nula y la hipótesis alternativa del test Dickey--Fuller, utilizado para evaluar la presencia de una raíz unitaria en el PIB real.

    \item[(d)] (8 puntos) A continuación se presentan los resultados de la estimación del modelo ADF descrito en el inciso (b):

    \begin{center}
        \begin{figure}[H]
            \centering
            \includegraphics[width=0.8\linewidth]{certamenes/Certamenes T2/rlst.png}
            \caption{Resultados del test ADF para el PIB real}
            \label{fig:adf_pib}
        \end{figure}
    \end{center}

   donde L(pib,1) representa el rezago de orden 1 del PIB. Utilizando estos resultados y los valores críticos apropiados del test Dickey--Fuller al 5\% de significancia, determine si existe raíz unitaria en el PIB real de Chile. 
\end{enumerate}
\newpage
    \item (20 puntos) Usando datos sobre gasto en transporte se estimó la siguiente ecuación para trabajadores de la región metropolitana


{\scriptsize
\begin{equation*}
\begin{array}{ll}
\widehat{GT} = & 1,2 \quad + \quad0,1 \cdot dist\quad + \quad 0,04 \cdot auto \quad - \quad 0,03 \cdot tele
 \\
& (0,5)  \hspace{0,8cm} (0,05) \hspace{1,3cm} (0,019) \hspace{1,3cm} (0,006)\\ 
& [0,5]  \hspace{0,9cm} [0,0511] \hspace{1,2cm} [0,02] \hspace{1,4cm}  [0,004] \\
& n = \text{450} \quad R^2 = \text{0,73}
\end{array}
\end{equation*}}



Aquí $GT$ es el gasto promedio en transporte en un mes que realiza un individuo; $dist$ es la distancia al trabajo; $auto$ cantidad de días que el individuo utiliza automóvil para ir al trabajo en el mes; $tele$ cuántos días de teletrabajo tiene el individuo en un mes. Entre paréntesis y entre corchetes se dan, respectivamente,
los errores estándar usuales y los errores estándar robustos a la heterocedasticidad.


\begin{enumerate}
    \item \textbf{(6 puntos)} ¿Tienen las variables  $dist$, $auto$ y $tele$ los efectos estimados esperados? ¿Cuáles de estas variables son estadísticamente significativas al nivel de 5\% suponiendo que se cumple el supuesto de homocedasticidad?

\item \textbf{(8 puntos)} Suponga que se sospecha heterocedasticidad, ¿se sigue manteniendo la significancia estadística al 5\%? Concluya respecto a lo obtenido.

\item \textbf{(4 puntos)} Plantee el test de Breusch-Pagan para heterocedasticidad. Sabiendo que el $R_{aux}^2$ de la regresión auxiliar alcanza el valor de 0.1, ¿es posible afirmar que hay heterocedasticidad?

\item \textbf{(4 puntos)} Se sospecha que puede haber multicolinealidad entre las variables, por lo que se corre la siguiente regresión:

$$auto = \delta_0 + \delta_1 \cdot dist - \delta_2 \cdot tele$$

Obteniendo un $R_j^2 = 0.89$. Con dicha información, ¿se puede afirmar que hay multicolinealidad?

\end{enumerate}
\end{enumerate}

\textbf{Pregunta extra}

Enuncie la condición de orden en contexto de ecuaciones simultáneas y dé una interpretación a cada término. (5 puntos extra)

\newpage
\section*{Formulario}
\begin{itemize}

\item Tabla $\chi^2$
\begin{center}
    \begin{tabular}{|c|c|c|c|}
    \hline
        df & 0.99 & 0.95 & 0.90 \\ 
        \hline
        1 &  6.635 & 3.841 & 2.706\\
        \hline
        2 &  9.210 & 5.991 & 4.605 \\
        \hline
        3 & 11.345 & 7.815 & 6.251  \\
        \hline
        4 &  13.277 & 9.488 & 7.779 \\
        \hline
        5 &  15.086 & 11.070 & 9.236 \\
        \hline
        6 &  16.812 & 12.592 & 10645 \\
        \hline
    \end{tabular}
\end{center}

\item Multiplicador de Lagrange

$$\color{red}{LM = nR^2 \sim \chi^2}$$

\item Test F para regresión auxiliar.
\begin{center}
 {  \color{red}     \(F=\dfrac{R_{\text{aux}}^{2}/q}{(1-R_{\text{aux}}^{2})/(n-k-1)}
          \sim F_{q,\;n-k-1}\).}
\end{center}



\item Tabla de test D-F
\begin{center}
    \begin{figure}[H]
        \centering
        \includegraphics[width=0.7\linewidth]{sem14/tabla.png}
    \end{figure}
\end{center}

\end{itemize}

\end{document}